\section{Approximation by truncating rare conditional excursions}
\label{sec:r76-rare-approximation}

Exact memory can be infinite although one finite simulator approximates
every stopped experiment with a prescribed expected cost.  The damping
construction of Theorem~\ref{thm:r75-stopped-approximation} gives such an
upper bound for the entire angle circle.  For a fixed finite command
alphabet we now obtain a smaller bound by preserving a long conditional
excursion exactly and modifying only its overflow.  The general mechanism
depends on the growth of the continuation semigroup, and applies beyond
the qubit family.

The ingredients have classical antecedents.  Polyhedral positive
realizations are part of the theory of Heller and of Monr\`as and Winter
\cite{Heller65,MonrasWinter14,MonrasWinter16}.  Martingale calculations
for pattern waiting times are classical \cite{Li80}.  The use made here
is a uniform approximation of complete adaptive stopped laws, with
an explicit memory bound from a finite return polytope and the number of
distinct continuation states.  It is separate from the exact noisy
realization problem of \cite{Fanizza24}.

\subsection{A sharp cost bound for a long conditional excursion}

Consider sequential calls, each producing a report declared either
\emph{continuing} or \emph{returning}.  Conditional on the entire past,
the chosen command and any private randomization, suppose the probability
of a continuing report is at most $r$, where $0<r<1$.
Let $\eta_L$ be the first call completing $L$ consecutive continuing
reports, with $L\ge1$.  A policy can stop initially or after any call;
its number of calls is $\tau$.

\begin{lemma}[Sharp expected-cost excursion bound]
\label{lem:r76-excursion-budget}
For every such policy with $\E\tau\le B<\infty$,
\begin{equation}\label{eq:r76-excursion-budget}
 \Prb\{\eta_L\le\tau\}
 \le \min\left\{1,\frac{(1-r)B}{r^{-L}-1}\right\}.
\end{equation}
The coefficient is optimal under the stated assumptions.  For independent
continuing indicators of probability $r$, the supremum over randomized
stopping policies of expected cost at most $B$ is exactly the right side.
\end{lemma}
\begin{proof}
Before the first completed run, let $K_t\in\{0,\ldots,L-1\}$ be the
current run length.  It increases by one on a continuing report and
returns to zero otherwise.  Absorb it at $L$ after the first completed
run, and put
\[
 u(k)=\frac{r^{-k}-1}{1-r}\quad(0\le k\le L).
\]
For $k<L$ and any conditional continuation probability $p\le r$,
\[
 p u(k+1)+(1-p)u(0)-u(k)
 \le r u(k+1)-u(k)=1.
\]
Stop this calculation at $\sigma_H=\min\{\tau,\eta_L,H\}$.
Summing the conditional inequalities, including the indicator that the
next call is actually made, gives
\[
 u(L)\Prb\{\eta_L\le\tau,\ \eta_L\le H\}
 \le \E u(K_{\sigma_H})\le\E\sigma_H\le B.
\]
Let $H\to\infty$.  Increasing events give
\eqref{eq:r76-excursion-budget}; no assumption about the expected duration
of another process is used.

For sharpness, take independent Bernoulli indicators of probability $r$.
The first completed run has finite expectation: disjoint blocks of length
$L$ succeed with probability $r^L$, so its expectation is at most
$L/r^L$.  In this case the displayed drift equals one.  Stopping at
$\eta_L\wedge H$ and then using bounded convergence for $u(K)$ and
monotone convergence for the duration proves
$\E\eta_L=u(L)$.  With probability
$\min\{1,B/u(L)\}$, selected before making any call, run this experiment
until $\eta_L$; otherwise stop immediately.  Its expected cost is at most
$B$ and its success probability equals the right side of
\eqref{eq:r76-excursion-budget}.
\end{proof}

Sharpness in this lemma concerns the probability of the exceptional
run.  It will not imply that every simulator must incur the approximation
error estimated by that probability.

\subsection{A general finite return polytope construction}

Let $\mathcal H$ be finite dimensional and let $\mathcal C$ be a closed convex set
of density matrices invariant under every positive-probability normalized
branch below.  The full state space is allowed.  For each command in a nonempty
finite set $A$, let $\{\mathcal I_{a,y}:y\in Y_a\}$ be a quantum
instrument, with finite $Y_a$.  Select a subset $R_a\subseteq Y_a$ of
continuing reports.  Suppose that there are density matrices
$\sigma_1,\ldots,\sigma_k$, with
$P=\operatorname{conv}\{\sigma_1,\ldots,\sigma_k\}$, such that
\begin{equation}\label{eq:r76-return-hypotheses}
 \begin{split}
 &\rho_0\in P,\qquad
   \sum_{y\in R_a}\tr\mathcal I_{a,y}(\rho)\le r<1,\\
 &\frac{\mathcal I_{a,y}(\rho)}{\tr\mathcal I_{a,y}(\rho)}\in P
 \quad(y\notin R_a,\ \tr\mathcal I_{a,y}(\rho)>0)
 \end{split}
\end{equation}
for every command $a$ and every $\rho\in\mathcal C$, with the seeds in $\mathcal C$.
On the full state space, the continuation bound is equivalently the finite
family of operator inequalities
$\sum_{y\in R_a}\mathcal I_{a,y}^*(I)\preceq rI$.
The return condition is a specified geometric hypothesis; it is not
inferred merely from mixing of the unobserved channel.  A prescribed
preparation in $\mathcal C$ but outside $P$ can be included by adding it to the seed list.

Write $\mathcal E=\{(a,y):a\in A,\ y\in R_a\}$ and $b=|\mathcal E|$.
For a word $w$ over $\mathcal E$, let $\mathcal I_w$ be its chronological
composition.  For $L\ge1$ define the finite set
\begin{equation}\label{eq:r76-orbit-dictionary}
 \mathscr D_L=
 \left\{\frac{\mathcal I_w(\sigma_j)}{\tr\mathcal I_w(\sigma_j)}:
  1\le j\le k,\ |w|<L,\ \tr\mathcal I_w(\sigma_j)>0\right\},
 \qquad G(L)=|\mathscr D_L|.
\end{equation}
The empty word is included.  Repeated states may be identified.

\begin{theorem}[Finite-memory approximation from return geometry]
\label{thm:r76-return-approximation}
Assume \eqref{eq:r76-return-hypotheses} with $0<r<1$.
For each $L\ge1$, there is a single $G(L)$-label mixed-start classical
model, chosen independently of the policy and of $B$, such that
\begin{equation}\label{eq:r76-return-approximation}
 \sup_{\pi:\,\E_{\rm target}^{\pi}\tau\le B}
 \|\mathcal L_{\rm target}^{\pi}-\mathcal L_L^{\pi}\|_{\rm TV}
 \le \min\left\{1,\frac{(1-r)B}{r^{-L}-1}\right\}.
\end{equation}
Here policies can choose commands, stopping and arbitrary terminal POVMs
from the visible record and their own randomness.  The complete stopped
record includes all commands and reports, the stopping time, the terminal
outcome and nontermination.  The bound also holds in trace distance for
the classical--quantum terminal record, with nontermination represented
by its scalar mass.  There is no additional entangled quantum register.

A counted pure start needs at most $G(L)+1$ labels.  In general
\begin{equation}\label{eq:r76-semigroup-growth}
 G(L)\le k\sum_{j=0}^{L-1}b^j.
\end{equation}
If the $b\ge1$ continuation maps commute pairwise, then
\begin{equation}\label{eq:r76-commuting-growth}
 G(L)\le k\binom{L+b-1}{b}.
\end{equation}
More generally any independently proved upper bound for the number of
distinct normalized states in \eqref{eq:r76-orbit-dictionary} can be used
in place of these bounds.  If that number is $O(L^d)$, the construction
uses $O((\log(B/\delta))^d)$ labels for fixed $r$ and
$B/\delta\to\infty$.
\end{theorem}
\begin{proof}
Use the states of $\mathscr D_L$ as hidden labels.  At a label $\rho$,
the probability of report $y$ on command $a$ is
$p_{a,y}(\rho)=\tr\mathcal I_{a,y}(\rho)$.  Zero probabilities require
no successor choice.  For a returning report, choose any convex
decomposition of its normalized output in the seed states.  For a
continuing report whose normalized output belongs to $\mathscr D_L$,
use that output label deterministically.  In the remaining case, send
the state to any one fixed seed.  Multiply each successor distribution
by $p_{a,y}(\rho)$ to obtain the report transition row.  All rows are
nonnegative and the sums over the reports are stochastic.  Use a seed
decomposition of $\rho_0$ for the initial distribution.  At label $\rho$,
the terminal decoder for any POVM $(E_z)_z$ is $\tr(E_z\rho)$.

To verify the error, first use all finite continuation words, without
the length restriction in \eqref{eq:r76-orbit-dictionary}.  This countable
refinement reproduces the target exactly.  Every returning branch is
expanded in the seeds; every continuing branch uses its exact normalized
image.  On each individual hidden label the weighted state identity is
exact, so induction gives the target subnormalized quantum state for
every command--report history.  Initial mixing, terminal POVMs and common
randomized policies preserve these identities.  This is only a proof
device, not an uncounted finite implementation resource.

Couple this exact refinement to the finite model by identical initial
seed choices and by identical reports and successor choices while the
states agree.  The coupling first fails only when a continuing image
leaves $\mathscr D_L$.  After a return the exact label is a seed, so a
product of fewer than $L$ continuing branches is in $\mathscr D_L$ by
definition.  The same is true initially.  Thus a failure requires a run
of at least $L$ consecutive continuing reports.  Identifying duplicate
states does not change this implication: their next normalized image
depends only on the state itself, and every shorter word image is still
present.

Before a failure the two visible histories, quantum states and policy
choices can be identical.  After a failure use arbitrary valid tails.
For a fixed finite horizon, the stopped classical--quantum states have
a common positive contribution from all paths with no failure.  If the
remaining mass is $e$, their trace distance is at most $e$, by positivity
and the triangle inequality.  The same assertion applies to the complete
record after any terminal POVM.  The target's probability of a failure
before stopping is at most $\Prb_{\rm target}^{\pi}\{\eta_L\le\tau\}$,
which is bounded by Lemma~\ref{lem:r76-excursion-budget}.  Its conditional
continuation bound follows from \eqref{eq:r76-return-hypotheses}, also
after conditioning on the virtual seed choices.

For unbounded stopping, retain every stopped history up to horizon $H$
and combine the unfinished histories into one symbol.  The sum of trace
norms on the finite stopped blocks increases to its full sum; the two
unfinished masses decrease to their respective nontermination masses.
Their absolute difference therefore converges to the nontermination
term.  The uniform finite-horizon bound passes to the limit, giving
\eqref{eq:r76-return-approximation}.  The budget is always evaluated under
the target.  The argument does not require the approximate model to have
expected duration at most $B$.

For a pure start append one initializer.  On the first actual command it
uses the initial mixture times the corresponding report matrix and
enters the core.  Its zero-call terminal decoder is that of $\rho_0$.
It is never entered again.  This proves the upper $G(L)+1$ without an
uncounted preliminary randomization.  If a separate finite policy uses
$s$ labels, the same construction gives at most $sG(L)+1$ labels for
the closed product implementation with a prescribed initial label.

There are $b^j$ words of length $j$, proving
\eqref{eq:r76-semigroup-growth}.  For commuting maps the product depends
only on a multi-index $m\in\mathbb N^b$ with $|m|<L$; there are
$\binom{L+b-1}{b}$ such multi-indices.  Normalization and deletion of zero
images cannot increase this number.  Finally it suffices to take
\[
 L\ge
 \max\left\{1,
 \left\lceil\frac{\log(1+(1-r)B/\delta)}{\log(1/r)}\right\rceil\right\}
 \quad(B>0,\ \delta>0).
\]
The asserted consequence of polynomial orbit growth follows.
\end{proof}

The construction preserves every finite history containing no run of
$L$ continuing reports, including its unnormalized terminal state.
It does not claim that its approximate response is stationary or that
the extra pure-start label is necessary.  When all continuation maps
have finite normalized orbit on the seeds, the full finite orbit instead
gives an exact model.  These are upper constructions, not a converse
classification of general positive realizations.

\subsection{A polylogarithmic simulator for finitely many commanded angles}

Apply the theorem to \eqref{eq:r75-noisy-instrument} with $\alpha=1$.
Let the finite known angle alphabet be
$\Theta=\{\theta_1,\ldots,\theta_a\}$, where $a\ge1$.
The continuation report is $0$, of probability $1/8$.
In the real Bloch plane put
\[
 v_0=(1,0),\qquad
 v_1=(-1/2,\sqrt3/2),\qquad
 v_2=(-1/2,-\sqrt3/2).
\]
Let $\sigma_j$ be the corresponding pure states.  Their triangle contains
the disk of radius $1/2$, and $I/2$ is their equal mixture.
Report $1$ returns to $\sigma_0=\tau$.  By
Lemma~\ref{lem:r75-full-instrument-bounds}, every normalized report-$2$
image has Bloch norm at most $5/11<1/2$, so it returns to this same
triangle, at every angle and from every real-plane state.
The real-plane density matrices form an invariant closed convex set:
the unitaries are real, and the other two branches preserve real symmetric
matrices.  Thus they supply the set $\mathcal C$ in the general theorem.

\begin{corollary}[Uniform finite-alphabet approximation]
\label{cor:r76-qubit-polylog}
For the undamped instrument at any fixed finite angle alphabet $\Theta$
and initial state $I/2$, there is, for every integer $L\ge1$, one classical
model with at most
\begin{equation}\label{eq:r76-qubit-count}
 3\binom{L+a-1}{a}
 \quad\hbox{mixed-start labels, or}\quad
 1+3\binom{L+a-1}{a}
 \quad\hbox{counted pure-start labels},
\end{equation}
whose complete stopped-law error satisfies
\begin{equation}\label{eq:r76-qubit-error}
 \sup_{\pi:\,\E_1^\pi\tau\le B}
 \|\mathcal L_1^\pi-\mathcal L_L^\pi\|_{\rm TV}
 \le \min\left\{1,\frac{7B}{8(8^L-1)}\right\}.
\end{equation}
The stronger stopped classical--quantum distance has the same bound.
The statements include adaptive choices among all supplied angles and
arbitrary terminal qubit POVMs.

For $B>0$ and $\delta>0$ it is enough to choose
\begin{equation}\label{eq:r76-qubit-depth}
 L=\max\left\{1,
   \left\lceil\log_8\left(1+\frac{7B}{8\delta}\right)\right\rceil\right\}.
\end{equation}
For a single commanded angle this uses at most $3L+1$ pure-start labels,
uniformly in that angle.  For each fixed finite alphabet it uses
$O_a((\log(B/\delta))^a)$ labels as $\delta/B\downarrow0$.
When the sine and cosine of each commanded angle are algebraic, every
state and transition coefficient of the displayed construction is
algebraic and is obtained by finitely many exact field operations.
\end{corollary}
\begin{proof}
The continuation maps are scalar multiples of commuting rotations, so
their normalized orbit dictionary consists of
\begin{equation}\label{eq:r76-qubit-dictionary}
 r_{j,m}=R_{m_1\theta_1+\cdots+m_a\theta_a}v_j,
 \qquad 0\le j<3,\quad m\in\mathbb N^a,\quad |m|<L.
\end{equation}
We may keep all these indices, even if some represent the same state.
Initialize uniformly on $(0,0),(1,0),(2,0)$.  From a label with Bloch
vector $r$, and command $\theta_i$, set
\[
 h=\tfrac18+\tfrac1{16}r_1,\qquad
 p_2=\tfrac78-h,\qquad
 u_2=-\tfrac18 R_{\theta_i}r-hv_0.
\]
Report $1$ has a row of mass $h$ into seed $(0,0)$.  Report $2$ has rows
into the three seeds of masses
\begin{equation}\label{eq:r76-explicit-return-rows}
 t_{2,j}=\frac{p_2+2v_j\mathbin{\cdot}u_2}{3},
 \qquad j=0,1,2.
\end{equation}
Indeed, $\sum_jv_j=0$ and
$\sum_jv_jv_j^{\mathsf T}=\tfrac32I_2$, so these masses sum to $p_2$
and their weighted Bloch vector is $u_2$.  Moreover,
$|u_2|/p_2\le5/11$ and $p_2\ge11/16$ give
\[
 t_{2,j}\ge\frac{p_2}{3}(1-10/11)\ge\frac1{48}>0.
\]
Thus the rows give the exact report-$2$ subnormalized state, with an
explicit positive margin.  On report $0$, whose mass is $1/8$, move
from $(j,m)$ to $(j,m+e_i)$ if $|m|<L-1$.  At $|m|=L-1$ send that mass
to the three seeds uniformly.  All report probabilities are unchanged
on each represented input state, and their sum is one.  Only the state
following an overflow zero report is modified.

The resulting table uses words with the same multi-index identified at
each depth and the same common-prefix proof as the theorem.  It keeps
the depth indices and may truncate at depth $L$ even when a periodic
phase makes that next physical state coincide with an earlier one.
The first possible modified state still occurs on the $L$th consecutive
zero.  Its state count
is \eqref{eq:r76-qubit-count}, and substitution of $r=1/8$ in
\eqref{eq:r76-return-approximation} gives
\eqref{eq:r76-qubit-error}.  Formula~\eqref{eq:r76-qubit-depth} makes
this bound at most $\delta$.  The label count and all error bounds are
uniform in the numerical values of the finite supplied angles, although
the transition table itself uses those values.  The explicit rotations
and \eqref{eq:r76-explicit-return-rows} prove the algebraic assertion.
\end{proof}

If $B=0$, a single initial decoder suffices, with no call and no error.
The damping upper in Theorem~\ref{thm:r75-stopped-approximation} remains
available and can be smaller for particular $B,\delta,a$; one may choose
the smaller of the two proved upper counts.  Its simultaneous
all-angle statement also remains valid.  In contrast, the finite dictionary
in \eqref{eq:r76-qubit-dictionary} is for a \emph{fixed finite} angle
alphabet; the bound is not asserted for a continuum of possible commands.

For every non-root phase, exact finite classical realization is still
impossible by Theorem~\ref{thm:r74-depolarizing-obstruction}.  At that same
fixed physical instrument the present bound is logarithmic in $B/\delta$
for a single angle.  Thus an inverse-square-root exact noise law is not a
lower bound on approximation of bounded-cost stopped experiments.
Section~\ref{sec:r76-physical-approximation} proves a matching logarithmic
law within the explicitly restricted class of physical dictionaries;
a matching lower for arbitrary classical terminal decoder tables remains
open.  The stopped-law bound also does not control the approximate
model's expected duration: unbounded costs
require a further tail or uniform-integrability hypothesis.  External
policy storage, table bits, arithmetic precision and exact sampling remain
separate resources.
